\subsection{LEFT DIFFERENTIAL}

DEFINITION 9. Let G be a Lie group, M a manifold of class \(\mathbf{C}^r\) and \(f\) a mapping of class \(\mathbf{C}^r\) of M into G. The left (resp. right) differential of \(f\) is the differential form of degree 1 on M with values in L(G) which associates with every vector \(u \in \mathrm{T}_m(\mathrm{M})\) the element \(f(m)^{-1}\) . \((\mathrm{T}_m f)(u)\) (resp. \((\mathrm{T}_m f)(u). f(m)^{-1}\) ).

In this chapter we shall only consider the left differential, which we shall denote by \(f^{-1}.df\) , and leave to the reader the task of translating the results for the right differential.

If \(f\) is the identity mapping of \(G, f^{-1}.df\) is the canonical left differential form \(\omega\) of \(G\) . Returning to the general case of Definition 8,

\[
(f ^ {- 1}. d f) _ {m} = \omega_ {f (m)} \circ \mathrm{T} _ {m} (f)
\]

and hence \(f^{-1}.df = f^{*}(\omega)\) . This implies that \(f^{-1}.df\) is of class \(\mathbf{C}^{r - 1}\) .

Examples. (1) If G is the additive group of a complete normable space and \(T_{0}(E)\) is canonically identified with E, \(f^{-1}\) . df is the differential df defined in Differentiable and Analytic Manifolds, R, 8.2.2.

\begin{enumerate}
\def\labelenumi{(\arabic{enumi})}
\setcounter{enumi}{1}
\tightlist
\item
  Suppose that G is the multiplicative group A* associated with a complete normable algebra A. Then f can be considered as a mapping of M into A and hence the differential df in the sense of Differentiable and Analytic Manifolds, R, 8.2.2 is defined and the product \(f^{-1}df\) in the sense of Differentiable and Analytic Manifolds, R, 8.3.2 is defined. Clearly the latter form is identical with the left differential of f.
\end{enumerate}

PROPOSITION 58. Let G and H be two Lie groups, M a manifold of class \(\mathbf{C}^r\) , \(f\) a mapping of class \(\mathbf{C}^r\) of M into G and \(h\) a morphism of G into H. Then

\[
(h \circ f) ^ {- 1}. d (h \circ f) = \mathrm{L} (h) \circ (f ^ {- 1}. d f) = (h ^ {- 1}. d h) \circ \mathrm{T} (f).
\]

For all \(x \in \mathbf{M}\) and \(u \in \mathrm{T}_x(\mathrm{M})\) ,

\[
(h \circ f) ^ {- 1}. d (h \circ f) (u) = ((h \circ f) (x)) ^ {- 1}. T (h \circ f) (u).
\]

The latter expression is equal, on the one hand, to

\[
\begin{array}{l l} \mathrm{T} (h) (f (x) ^ {- 1}. \mathrm{T} (f) (u)) & (\S 2, \text { Proposition   5 }) \\ = \mathrm{T} _ {e} (h) ((f ^ {- 1}. d f) (u)) \end{array}
\]

and, on the other hand, to

\[
\begin{array}{c} h (f (x)) ^ {- 1} \mathrm{T} (h) (\mathrm{T} (f) u) \\ = (h ^ {- 1}. d h) (\mathrm{T} (f) u). \end{array}
\]

PROPOSITION 59. Let G be a Lie group, M a manifold of class \(\mathbf{C}^r\) , \(f\) and \(g\) mappings of class \(\mathbf{C}^r\) of M into G and \(p\) the canonical surjection of TM onto M.

\[
(i) (f g) ^ {- 1}. d (f g) = (\mathrm{Ad} \circ g \circ p) ^ {- 1} \circ (f ^ {- 1}. d f) + g ^ {- 1}. d g.
\]

\begin{enumerate}
\def\labelenumi{(\roman{enumi})}
\setcounter{enumi}{1}
\tightlist
\item
  Writing \(h(m) = f(m)^{-1}\) for all \(m \in M\) ,
\end{enumerate}

\[
h ^ {- 1}. d h = - (\mathrm{Ad} \circ f \circ p) \circ (f ^ {- 1}. d f).
\]

Assertion (i) follows from § 2, no. 2, Proposition 7. Assertion (ii) follows from (i) by putting \(g = h\) .

COROLLARY 1. Let \(s \in G\) and \(sg\) be the mapping \(x \mapsto sg(x)\) of \(M\) into \(G\) . Then \((sg)^{-1}.d(sg) = g^{-1}.dg\) .

This follows from Proposition 59 (i) taking \(f\) to be the constant mapping \(x \mapsto s\) of \(M\) into \(G\) .

COROLLARY 2. If the mappings \(f\) and \(g\) of \(M\) into \(G\) have the same left differential, the tangent mapping to \(fg^{-1}\) is everywhere zero. If further \(K\) is of characteristic 0, then \(fg^{-1}\) is locally constant.

By Proposition 59,

\[
(f g ^ {- 1}) ^ {- 1}. d (f g ^ {- 1}) = (\mathrm{Ad} \circ g \circ p) \circ (f ^ {- 1}. d f) - (\mathrm{Ad} \circ g \circ p) \circ (g ^ {- 1}. d g).
\]

If \(f^{-1}.df = g^{-1}.dg\) , then \((fg^{-1})^{-1}.d(fg^{-1}) = 0\) , that is \(T_x(fg^{-1}) = 0\) for all \(x \in M\) . This proves the first assertion. The second follows from it by Differentiable and Analytic Manifolds, R, 5.5.3.

PROPOSITION 60. Let G be a Lie group, of finite dimension if K is of characteristic \textgreater0, M a manifold of class \(\mathbf{C}^r\) , \(f\) a mapping of class \(\mathbf{C}^r\) of M into G and \(\alpha\) the left differential of \(f\) . Then \(d\alpha + [\alpha]^2 = 0\) .

Let \(\omega\) be the canonical left differential form of G. Using Corollary 1 to Proposition 51, no. 14, we have

\[
\begin{array}{r l} d \alpha & = d (f ^ {*} (\omega)) = f ^ {*} (d \omega) = f ^ {*} (- [ \omega ] ^ {2}) \\ & = - [ f ^ {*} (\omega) ] ^ {2} = - [ \alpha ] ^ {2}. \end{array}
\]

